\chapter{Trace the decoder's remembered state}
\MARSChapterHeading{Trace the decoder's remembered state}
\label{revision-lab:qoi}

A compressed pixel does not always carry its own components. It may repeat the
previous color, recover a cached color, or describe a difference from the previous
one. To explain a decoded pixel, you therefore need both the selected source
bytes and the state established by earlier operations. This lab follows that
state through a complete decoder before asking you to change its output placement.

Import the \texttt{030-qoi} student archive and open it in Learn. The entry
\texttt{qoi-codec.asm} selects the complete eight-pixel decoder. The
\texttt{worked/} copies select other callers or change one input condition.
The independent \texttt{qoi\_decode\_mirrored} entry is deliberately unfinished.
Read the main chapter's format, cache and publication rules first; use this
companion to connect the observations with the reasons for each boundary.

Answer predictions inline. Each question's source pointer refers to its selected
example. During Debug, inspect the selected paused frame: after Learn recognizes
a pause, its outline can move to the next requested instruction. A highlighted
instruction has not executed yet. Labels beginning \texttt{after\_} mark
checkpoints after the preceding operations. Use one breakpoint at a time,
resume between observations and finish a process before starting another source.
Keep these supplied examples unchanged; use disposable copies for further faults.

\section{Keep four different lengths separate}

An eight-by-one image requires eight output pixels, or 32 raw RGBA bytes.
Its selected QOI input contains 38 bytes. A conservative encoding reservation
for those dimensions is $22+5\cdot8=62$ bytes: fourteen header bytes, eight
marker bytes, and at most five bytes per pixel. These quantities answer different
questions. A reservation says what could fit; it does not say which bytes have
actually arrived or how many pixels a particular parse has reconstructed.

\Needspace{4\baselineskip}
The bound observation uses \texttt{worked/bound.asm}. At
\texttt{after\_b\_size}, the arithmetic is complete: status is zero and
\texttt{\$v1} is 62. The same result reaches \texttt{bound\_return}.
The caller has not yet stored its result slots. Continue to termination only
after retaining the returned value; the exit service later reuses \texttt{\$v0}.

Return to the decode entry for the header observation. Its serialized width is
\texttt{00 00 00 08}, and its height is \texttt{00 00 00 01}. Unsigned byte
loads, shifts and OR reconstruct eight and one in \texttt{\$t4} and
\texttt{\$t5}. A native little-endian word load would give a different value.
At this boundary \texttt{\$a1} is still the selected encoded length, 38.
The routine has not converted that argument into a pixel count.

The outer decoder first calls the validator. At \texttt{after\_d\_call},
validation returns eight required pixels in \texttt{\$v1}. Later, at
\texttt{decode\_return}, decode returns packed dimensions
\texttt{0x0000000100000008}. That value is height one in the upper word and
width eight in the lower word. It is neither 38 encoded bytes nor 32 output bytes.
For a descriptor-only observation, the complete \texttt{worked/validate.asm}
copy publishes native words eight, one, four and zero at \texttt{validate\_return}.

\section{Follow one transition before following the whole image}

Before the first opcode, the previous-color predictor is opaque black
$(0,0,0,255)$, while every cache slot is transparent black $(0,0,0,0)$.
At the first \texttt{after\_w\_clear}, \texttt{\$s3} is
\texttt{00000000ff000000}, \texttt{\$s2} is eight, and
\texttt{\$s6} is the sentinel minus one for no previous INDEX. The walker is
validating, so \texttt{\$fp} is zero and caller pixels are not emitted.

The first chunk is \texttt{c0}. Its low six bits are zero, but RUN stores
count minus one. It therefore repeats the predictor once. At
\texttt{after\_w\_run}, one repetition has been decoded into \texttt{\$t3};
eight pixels remain because the advance has not happened. Reading a count and
consuming its output are separate transitions.

\begin{table}[htbp]
\centering
\begin{tabular}{l r r r}
\toprule
Chunk bytes & Bytes consumed & Pixels produced & Pixels left\\
\midrule
\texttt{c0} & 1 & 1 & 7\\
\texttt{fe ff 00 00} & 4 & 1 & 6\\
\texttt{c0} & 1 & 1 & 5\\
\texttt{35} & 1 & 1 & 4\\
\texttt{7f} & 1 & 1 & 3\\
\texttt{a5 97} & 2 & 1 & 2\\
\texttt{ff 07 06 05 80} & 5 & 1 & 1\\
\texttt{04} & 1 & 1 & 0\\
\bottomrule
\end{tabular}
\caption{Account for the worked stream after each operation's output has
advanced. The chunks consume sixteen bytes and produce eight pixels. Header and
marker add 22 bytes, giving the selected file length 38. A longer RUN would
produce several pixels while consuming only one source byte.}
\label{tab:revision-qoi-lab-counts}
\end{table}

For each operation, retain two independent accounting statements: the cursor
has consumed its own encoded bytes, and the remaining pixel count has decreased
by its produced pixels. A cursor at the marker does not alone establish the
right number of pixels. A zero remaining count does not authorize extra payload
or a damaged marker. Both conditions must hold.

\section{Explain cache changes and complete colors}

The first RUN installs its opaque-black color in decimal cache slot 53.
At \texttt{after\_w\_hash}, \texttt{\$s4} identifies the private cache base;
slot 53 is $53\cdot4=212$ bytes from it. Use that live base if you inspect Memory.
A validation-only caller and a decoder that calls validation have different
nested frames, so an address printed in an example is not a fixed cache address.

\begin{figure}[htbp]
  \centering
  \begin{tikzpicture}[alt={Actual light-theme Memory capture during the first validation RUN. The row begins at private cache slot 52; slot 53 contains bytes 00 00 00 ff while neighboring slots remain zero. A native box identifies the inserted word.}]
    \node[anchor=north west,inner sep=0,font=\small\bfseries] (title) at (0,0) {RUN inserts opaque black into slot 53};
    \node[anchor=north west,inner sep=0] (a) at ([yshift=-2mm]title.south west) {\MARSRevisionCrop{revision/captures/qoi/qoi-run-cache.png}{0.99\linewidth}{0.03125}{0.31555556}{0.58203125}{0.60888889}};
    \begin{scope}[shift={(a.south west)},x={($(a.south east)-(a.south west)$)},y={($(a.north west)-(a.south west)$)}]
      \draw[black,line width=.8pt] (.39,.3088) rectangle (.58,.6176);
    \end{scope}
  \end{tikzpicture}
  \caption{The row begins at slot 52. Slot 53 contains opaque black after the first RUN; nearby zero entries remain transparent black. These private addresses describe the supplied walker, not a required layout for equivalent solutions. Native LaTeX supplies the title, crop and outlines.}
  \label{fig:revision-qoi-native-cache}
\end{figure}


The inserted bytes are \texttt{00 00 00 ff}. Nearby cleared slots still contain
\texttt{00 00 00 00}. RGB black is the same, but alpha is different. The
fourth pixel's INDEX byte \texttt{35} selects slot 53 and must recover opaque
black. A validator that tracks only opcode widths would miss this dependency.
It must reconstruct the same predictor and cache state as emission.

In the chunk observation, \texttt{after\_w\_index} occurs twice. On the first
requested visit, \texttt{\$s3} is \texttt{ff000000} and five pixels remain
before advancing the fourth pixel. On the later visit, \texttt{\$s3} is
\texttt{ff010101} and one remains before advancing the eighth. The source label
is the same; the incoming history explains why the states differ. Resume through
the intervening operations rather than restarting when Learn requests that label
again.

For DIFF \texttt{7f}, each two-bit field is three, with bias two, so each RGB
component increases by one. The preceding black becomes $(1,1,1)$.
For LUMA \texttt{a5 97}, $d_g=37-32=5$, $d_r-d_g=9-8=1$ and
$d_b-d_g=7-8=-1$. Starting from $(1,1,1)$ gives $(7,6,5)$.
At \texttt{after\_w\_diff} and \texttt{after\_w\_luma}, the components
are still separate temporaries; the following pack constructs the predictor word.
Apply modulo-256 reduction to each component. Clamping would decode another format.

The next RGBA literal supplies $(7,6,5,128)$. The last INDEX restores
$(1,1,1,255)$, including its earlier alpha. An INDEX does not retain alpha from
the immediately preceding predictor. The complete decoded words are visible below.

\begin{figure}[htbp]
  \centering
  \begin{tikzpicture}[alt={Actual light-theme Memory capture after decoding the eight-pixel worked stream. The two displayed rows contain all eight framebuffer words. The seventh pixel has byte sequence 05 06 07 80 and the eighth restores alpha ff. A native box marks the seventh word.}]
    \node[anchor=north west,inner sep=0,font=\small\bfseries] (title) at (0,0) {Eight pixels with complete alpha};
    \node[anchor=north west,inner sep=0] (a) at ([yshift=-2mm]title.south west) {\MARSRevisionCrop{revision/captures/qoi/qoi-trace-rgba.png}{0.99\linewidth}{0.03125}{0.31555556}{0.58203125}{0.60888889}};
    \begin{scope}[shift={(a.south west)},x={($(a.south east)-(a.south west)$)},y={($(a.north west)-(a.south west)$)}]
      \draw[black,line width=.8pt] (.58,.005) rectangle (.77,.303);
    \end{scope}
  \end{tikzpicture}
  \caption{All eight reconstructed pixels appear as native little-endian word bytes. The outlined word retains alpha hexadecimal 80; the next word restores alpha ff. Native LaTeX supplies the title, crop and outlines.}
  \label{fig:revision-qoi-native-rgba}
\end{figure}


The seventh word \texttt{80070605} appears in byte view as
\texttt{05 06 07 80}, because native words are little-endian. It still means
$(R,G,B,A)=(7,6,5,128)$. The descriptive-flags copy changes channels to three
and colorspace to one. Its RGBA literal and result retain those same components;
neither flag asks this codec to discard alpha or transform the bytes. An RGB
preview cannot establish this property because it omits the fourth component.

\section{Start each parse with the same initial state}

Validation updates private state while proving the complete input. Once it
returns, the outer decoder compares the required count with the saved unsigned
capacity. At \texttt{after\_d\_capacity}, eight required pixels fit 4096.
Only then does the second walk start emitting. At \texttt{after\_d\_walk},
the caller pixels exist and the walker has returned eight pixels. The outer
routine still has to construct its packed-dimension return value.

Notice that \texttt{\$s3} means saved capacity at the outer decoder's boundary,
but predictor color inside the walker. Saved registers are preserved across
calls; they are not permanently assigned one meaning across all routines.
Identify the current routine before interpreting a displayed value.

Restarting emission needs more than the original source pointer. Its predictor,
all 64 cache entries and prior-INDEX history must start fresh. Two small examples
make errors here visible without relying on the familiar trace.

The initial-INDEX copy begins with INDEX zero. An untouched cache slot is defined
transparent black, so this valid one-pixel image produces word
\texttt{00000000}. It does not use the opaque initial predictor and does not
require a previous literal. Inspect the first validation INDEX, then the complete
one-by-one decode return and framebuffer.

\Needspace{12\baselineskip}
The cache-reset copy begins with the same INDEX, then supplies RGBA $(1,0,0,23)$.
That late color hashes to
\[
 (3\cdot1+5\cdot0+7\cdot0+11\cdot23)\bmod64
 =256\bmod64=0.
\]
Validation therefore ends with slot zero populated. At \texttt{after\_d\_call},
the input has been accepted but no caller pixel has been emitted. The emitting
pass's first INDEX must nevertheless recover zero. If it inherits validation's
cache, its first pixel changes to the \emph{later} color. The expected output
words are \texttt{00000000} followed by \texttt{17010000}. Explain why a trace
that overwrites each slot before reading it could conceal that error.

\section{Reject the whole input before publishing a prefix}

The late-marker copy keeps all eight plausible pixel operations but replaces
the final marker one with zero. It also gives the decoder zero pixel capacity.
Predict the precedence before launching it. At \texttt{decode\_return}, expect
format status one and zero dimensions. Validation fails before capacity is tested,
and the complete destination retains its initial \texttt{a5} bytes.

\begin{figure}[htbp]
\centering
\begin{tikzpicture}[alt={Actual light-theme Memory capture at decode return after the final marker byte was changed from one to zero. The displayed framebuffer row remains sixteen a5 sentinel bytes. Separate executable checks establish the whole-arena failure guarantee.}]
\node[anchor=north west,inner sep=0,font=\small\bfseries] (title) at (0,0) {The first row remains untouched};
\node[anchor=north west,inner sep=0] (a) at ([yshift=-2mm]title.south west) {\MARSRevisionCrop{revision/captures/qoi/qoi-rejected-frame.png}{0.99\linewidth}{0.03125}{0.31555556}{0.58203125}{0.60888889}};
\begin{scope}[shift={(a.south west)},x={($(a.south east)-(a.south west)$)},y={($(a.north west)-(a.south west)$)}]
\draw[black,line width=.7pt] (.20,.3088) rectangle (.965,.6176);
\end{scope}
\end{tikzpicture}
\caption{The displayed first row remains unchanged after a late marker error. Whole-arena and access checks separately verify absence of rejected output writes. Native LaTeX supplies the title, crop and outlines.}
\label{fig:revision-qoi-native-rejection}
\end{figure}


The visible row helps locate the output, but it is not evidence for the whole
arena. Executable checks inspect the complete destination and its access history.
The valid short-capacity copy separates the other failure: seven available pixels
cannot hold eight required ones. At \texttt{decode\_capacity}, success and the
required count are still live before the highlighted \texttt{li} constructs
status two. At the return, the count is zero and no pixels have been published.

For a truncated RGBA operation, a missing alpha byte may be physically readable
as a marker zero. That does not make it an operand. Draw the selected file end
and the boundary eight bytes before it; prove availability against the chunk
boundary before loading components. Keep this source permission separate from
the later proof that enough destination pixels exist.

\section{Follow the initialized file prefix}

The short-read activity supplies virtual \texttt{input.qoi} containing the same
38 bytes as \texttt{worked/input.qoi}. Positive reads transfer at most seven
bytes, and the first acquired descriptor is 19. These are this activity's supplied
conditions; ordinary execution outside it uses the real file and normal file
services. A fresh launch starts independent input state.

At \texttt{after\_l\_open}, the loader owns descriptor 19 and has accepted zero
bytes. At the first \texttt{after\_l\_read}, both the returned transfer count
and initialized prefix are seven. The large remaining-capacity request did not
initialize the rest of the buffer. Continue accumulating actual positive results;
a short positive result is not EOF. The final positive transfer is three bytes,
followed by a zero result that establishes the end of this 38-byte file.

At \texttt{load\_return}, the loader has closed its descriptor and validated
the complete file. Its result length is 38 encoded bytes. Compare that with the
validator's eight-pixel result and decoder's packed dimensions. Finish the silent
ordinary Run of the decode entry as well. Its empty console establishes normal
termination, while the earlier Debug and Memory observations explain the numerical
result. The next experiment uses actual host files and a complete pipeline.

\section{Produce actual files and compare their decoded meaning}

From the public book repository, run the complete pipeline in a new directory,
substituting Java and CLI paths:
\begin{verbatim}
python3 tools/run_revision_qoi.py \
  --java /path/to/java --cli /path/to/cli.jar \
  --output /tmp/qoi-trace
\end{verbatim}
The helper writes a deterministic \texttt{input.qoi}, selects mode six in a
separate source copy and executes ordinary host-file services. It checks the
output's strict local format, dimensions, descriptive flags and complete RGBA.
It also begins with a longer old output to check successful truncation.
Each directory retains source, input, output, raw decoded RGBA and a receipt.
It refuses an existing output directory.

The worked encoder produces 35 bytes from the 38-byte trace. Compare the first
red chunk: input \texttt{fe ff 00 00} occupies four bytes, while output
\texttt{5a} represents differences $(-1,0,0)$ from black. Wrapping red's
$0-1$ gives 255. The three-byte saving changes the representation without
changing a pixel. Equivalent encoders may make other legal choices.

Repeat with \texttt{--fixture sprite}. The resulting image reconstructs the
familiar eight-by-eight pattern. Its 256 raw RGBA bytes are represented by a
61-byte file in this supplied example; that is a byte-size observation, not a
speed measurement. The actual native Bitmap observation below shows its RGB.

\Needspace{22\baselineskip}
\begin{figure}[!ht]
  \centering
  \begin{tikzpicture}[alt={Actual native light-theme Bitmap Display of the decoded eight-by-eight sprite. White cells correspond to the set bits of rows hexadecimal 16,3c,7e,db,ff,24,42,81 against a black background. All sixty-four RGB cell interiors match the independently specified fixture. Native LaTeX identifies the dimensions and frames the original pixel canvas.}]
    \node[anchor=north west,inner sep=0] (a) {\MARSRevisionCrop{revision/captures/qoi/qoi-decoded-sprite.png}{0.48\linewidth}{0.20390625}{0.06666667}{0.20390625}{0.09111111}};
    \node[anchor=south west,inner sep=0,font=\small\bfseries] at ([yshift=2mm]a.north west) {Eight columns by eight rows};
    \draw[black,line width=.7pt] (a.south west) rectangle (a.north east);
  \end{tikzpicture}
  \caption{Actual native Bitmap rendering after decoding the sprite. The capture checks all 64 RGB cell interiors; separate word and host-decoder checks verify complete RGBA. Native LaTeX supplies the crop, dimension label and frame.}
  \label{fig:revision-qoi-native-sprite}
\end{figure}


To reproduce the Bitmap view, use the helper's sprite pipeline copy and break
at \texttt{decode\_return}. Request the same view with these public source comments; settings for the same
tool are combined across the two lines:
\begin{verbatim}
# mars:tool bitmap width=256 height=256 unitWidth=32 unitHeight=32
# mars:tool bitmap base=custom address=0x10015180
\end{verbatim}
This address is the supplied framebuffer's ordinary layout; use its actual
address if you relocate objects. The header configures the viewer, not the codec.

The helper also writes \texttt{rgb-preview.ppm}. That preview deliberately omits
alpha, while \texttt{decoded.rgba} retains it. The draft verification separately
decodes output files with a pinned build of the official QOI implementation and
compares all four components and flags. The preview alone is not that evidence.

Use \texttt{--fixture pattern --width 64 --height 64} for the largest local
image. Its generated input consists of RGBA literals and occupies exactly
20502 bytes, exercising the full-buffer EOF boundary. Try
\texttt{--channels 3 --colorspace 1} with the trace to keep the alpha observation
separate from descriptive flags. These ordinary host runs do not force short
transfers or injected errors; grading fixtures test those paths separately.

\section{Place reflected pixels while keeping forward decoder state}

Complete \texttt{qoi\_decode\_mirrored} in the delivered
\texttt{independent/qoi-mirrored.asm}. Keep the supplied caller and
helpers. The interface matches decode: byte pointer, unsigned byte count,
word destination and unsigned pixel capacity. Preserve status precedence,
zero failure length, full saved state and all source bytes. Format or capacity
failure must leave the entire destination untouched.

The output maps $(x,y)$ to $(W-1-x,y)$ and returns the original packed dimensions.
One-column images remain unchanged; odd-width rows retain their center pixel.
Reverse each row, not the whole framebuffer. A run may cross rows in the stream,
so a decoder that emits directly to reflected addresses still needs to compute
placement for every produced pixel.

\begin{figure}[htbp]
\centering
\begin{tikzpicture}[alt={Actual light-theme Memory capture of the independent three-by-two reflected result. Six words preserve their alpha while red values appear in order three, two, one, six, five, four. The first image row occupies twelve bytes, whereas one Memory row displays sixteen. Unused words remain a5 sentinels.}]
\node[anchor=north west,inner sep=0,font=\small\bfseries] (title) at (0,0) {Image rows and Memory rows have different widths};
\node[anchor=north west,inner sep=0] (a) at ([yshift=-2mm]title.south west) {\MARSRevisionCrop{revision/captures/qoi/qoi-mirrored-rows.png}{0.99\linewidth}{0.03125}{0.31555556}{0.58203125}{0.60888889}};
\begin{scope}[shift={(a.south west)},x={($(a.south east)-(a.south west)$)},y={($(a.north west)-(a.south west)$)}]
\draw[black,line width=.8pt] (.20,.3088) rectangle (.77,.6176);\draw[black,line width=.8pt,dashed] (.77,.3088) rectangle (.965,.6176);\draw[black,line width=.8pt,dashed] (.20,.005) rectangle (.58,.303);
\end{scope}
\end{tikzpicture}
\caption{Six pixels preserve their full words while each three-pixel image row is reflected. The solid outline marks the first image row; the dashed outlines mark the second image row across two Memory rows. Native LaTeX supplies the title, crop and outlines.}
\label{fig:revision-qoi-native-mirror}
\end{figure}


This observation uses a three-by-two fixture. The original red values are
one, two, three in the first row and four, five, six in the second, with varied
alpha. The reflected rows are three, two, one and six, five, four. A Memory row
shows four words, while an image row here has only three. Thus the first word of
the second image row appears at the end of the first displayed Memory row.
The unused bytes after six pixels retain their sentinel pattern.

You may decode successfully into the destination and then swap its row words,
or use a forward state walk with reflected stores after complete validation.
Explain where your design proves the whole input and where it checks capacity.
Changing an emission address must not change which pixel becomes the predictor
or which color is inserted into the cache.

\begin{verbatim}
java -jar /path/to/cli.jar grade \
  revision/labs/checks/qoi-mirrored.kts \
  /path/to/your/qoi-mirrored.asm
\end{verbatim}
The untouched starter deliberately stops in the unfinished entry. Learn supplies
seven check groups: dimensions/capacities, chunk/flag behavior, alignment, malformed
files, rejected full-width counts, complete caller execution and initial-state
counterexamples. Together they retain the original 205 finite cases and budgets.
Use the first failure to choose a concrete observation; graduated hints identify
which permission or state argument to revisit. The optional CLI route above checks
the same independent task. Both accept equivalent implementations. They do not establish every possible
QOI stream. Test a successful call after a failure and retain an explanation of
which state must start fresh.

Finally run your completed routine through the real file pipeline:
\begin{verbatim}
python3 tools/run_revision_qoi.py \
  --java /path/to/java --cli /path/to/cli.jar \
  --source /path/to/your/qoi-mirrored.asm --mirror \
  --fixture pattern --output /tmp/qoi-mirrored
\end{verbatim}
The helper selects your entry in a separate pipeline copy and compares the
encoded output's reconstructed image with the expected reflection. Keep the
artifact, receipt and a diagram of one multirow placement. Those explain more
than a claim that a symmetric sprite still looks correct.
